\chapter{실제 LLM 행동에서 발견된 메커니즘}
\label{chap:phase7}

\section{결론부터: 발견은 구체적이며 보편성은 미확립이다}

현재 문헌은 자연 학습 LLM의 일부 행동에서 재현 가능한 내부 구조를 발견했지만, 그 결과는 대체로 특정 모델, prompt distribution, token position과 metric에 한정된다.
induction, IOI와 사실 회상은 중요한 실증 사례이지만 모든 자연어 처리에 공통인 하나의 전역 알고리즘으로 합쳐지지 않았다.

행동별 연구를 읽을 때는 ``무엇을 발견했는가''와 ``어디까지 발견했는가''를 같은 문단에 기록해야 한다.
이 장은 각 결과 뒤에 모델, 과제, 개입과 남은 대안을 붙이며, 사람에게 익숙한 이야기만으로 빈 부분을 메우지 않는다.

\section{factual recall}

사실 회상 연구의 가장 통제된 형태는 subject--relation--attribute triple에서 다음 token을 예측하는 과제다.
예를 들어
\begin{equation}
  (s,r,a)=(\text{Barack Obama},\text{place of birth},\text{Honolulu})
\end{equation}
를 ``Barack Obama was born in the city of''라는 prompt로 표현하고, 모델이 attribute의 첫 token을 예측하게 한다.
이 설정은 open-ended generation보다 target logit과 token position을 명확히 정의한다.
\index{factual recall@사실 회상(factual recall)}

이 과제의 정확성은 모델이 세계 사실을 추론하는 모든 방식을 포괄하지 않는다.
attribute가 여러 token이면 연구가 첫 token만 설명할 수 있고, prompt paraphrase나 여러 정답이 있는 relation에서는 metric이 달라진다.
학습 corpus에 사실이 어떤 문맥과 빈도로 있었는지도 대개 완전히 알 수 없다.

ROME의 causal tracing은 GPT-2 계열에서 subject token을 처리하는 middle-layer MLP가 factual prediction에 중요한 역할을 한다는 증거를 제시했고, rank-one weight update로 선택한 factual association을 편집했다 \citep{meng2022rome}.
그러나 편집 가능성과 지식의 유일한 저장 위치는 같은 명제가 아니다.
후속 분석은 localization 방법과 editing 성능이 항상 일치하지 않는다고 보고했다 \citep{hase2023localization}.

\section{factual recall의 세 단계: 원 연구의 정확한 주장}

Geva 등의 원 연구가 제안한 세 단계는 ``subject enrichment, relation propagation, attribute extraction''이다 \citep{geva2023recall}.
연구 대상은 정답 attribute를 다음 token으로 예측한 GPT-2 XL 1.5B와 GPT-J 6B의 각각 약 1,200개 query이며, attention knockout, vocabulary projection과 component cancellation을 사용했다.

첫 단계는 early MLP가 마지막 subject position의 표현을 여러 subject-related attribute로 풍부하게 만든다는 결과다.
연구는 last-subject representation을 vocabulary에 투영한 상위 token에서 Wikipedia로 구성한 subject attribute 집합의 비율을 측정했다.
GPT-2에서 early MLP update를 취소하면 layer 40의 attribute rate가 평균 약 88\% 감소했고, MHSA update를 취소한 경우에는 감소가 30\% 미만이었다고 보고했다.

둘째 단계는 relation 정보가 subject 정보보다 먼저 prediction position으로 전파된다는 결과다.
연구는 last position이 non-subject position을 읽는 attention edge를 연속 layer window에서 막았을 때 나타나는 정답 확률 변화를 측정했다.
그 영향 peak가 subject-position edge를 막았을 때의 peak보다 먼저 나타났기 때문에 relation propagation으로 해석했다.

셋째 단계는 upper-layer MHSA가 relation을 담은 prediction representation으로 enriched subject를 조회하여 target attribute를 추출한다는 결과다.
연구는 last position이 subject position을 읽는 attention을 분석하고, 일부 head의 parameter가 subject--attribute mapping을 지원한다고 보고했다.
저자들은 조사한 prediction의 약 70\%에서 이 extraction behavior를 관찰했다고 명시했다.

\begin{figure}[htbp]
\centering
\begin{tikzpicture}[
  node distance=8mm and 10mm,
  stage/.style={draw=Navy,rounded corners,fill=SoftBlue,text width=38mm,
                minimum height=17mm,align=center,font=\small},
  arr/.style={-{Latex[length=2.2mm]},thick,Navy}
]
  \node[stage] (s1) {1. subject enrichment\\마지막 subject position\\early MLP 중심};
  \node[stage,right=of s1] (s2) {2. relation propagation\\relation 정보가\\last position으로 이동};
  \node[stage,right=of s2] (s3) {3. attribute extraction\\upper MHSA가 subject에서\\관계에 맞는 속성을 추출};
  \draw[arr] (s1) -- (s2);
  \draw[arr] (s2) -- (s3);
\end{tikzpicture}
\caption{Geva et al. (2023)이 제안한 사실 회상의 세 단계}
\label{fig:geva-three-stages}
\end{figure}

이 세 단계는 software의 세 개 함수처럼 완전히 분리된 시간 구간이 아니다.
layer window별 개입 효과와 projection 결과를 세 가지 기능으로 요약한 경험적 메커니즘이며, 여러 component의 계산은 겹칠 수 있다.
``모델이 먼저 모든 속성을 명시적 목록으로 만들고 다음에 데이터베이스 검색을 실행한다''는 literal한 절차도 아니다.

\subsection{2510.09033의 Section 3.3은 무엇을 다시 정의하는가}

Cheang 등은 원 연구의 세 단계를 hallucination 분류를 위한 세 intervention component로 다시 기술한다 \citep{cheang2026knowledge}.
그 논문의 표현은 early-layer subject enrichment, middle-layer subject-to-last attention flow, upper-layer last-token decoding이다.
특히 둘째 항목을 ``subject information propagation''으로 부르므로, Geva 원 연구의 ``relation information이 먼저 last position으로 전파된다''는 둘째 단계와 문구가 완전히 같지 않다.

Section 3.3의 문장은 token 생성 algorithm을 새로 제안하는 부분이 아니라 label을 만드는 근거를 설명하는 부분이다.
저자들은 오답을 낸 표본에서 last position이 subject token을 읽는 attention을 막고, 원 출력분포와 개입 후 출력분포의 Jensen--Shannon divergence를 계산한다.
이 값이 factual association 표본 전체의 평균보다 크면 associated hallucination, 그렇지 않으면 unassociated hallucination으로 분류한다.

\begin{factbox}
``세 단계를 거쳐 completion을 produce한다''는 말은 token 하나가 세 번 출력된다는 뜻이 아니다.
한 번의 forward pass 안에서 서로 다른 layer와 position의 activation이 subject를 풍부하게 만들고, 필요한 정보를 마지막 position으로 전달하며, 최종 logit에서 attribute token을 높이는 기능적 분해를 뜻한다.
\end{factbox}

\begin{limitbox}
2510.09033의 FA/AH/UH는 intervention과 threshold로 만든 조작적 범주다.
모든 hallucination이 자연적으로 이 세 ontological class 중 하나를 가진다는 증명은 아니다.
특히 AH/UH label이 subject-reliance intervention으로 정의되므로, 그 뒤의 subject-flow 차이는 독립적인 진실성 label만 분석한 결과로 읽어서는 안 된다.
\end{limitbox}

\section{hallucination과 내부 상태}

Cheang 등의 핵심 결과는 조사한 두 모델에서 associated hallucination의 내부 상태가 factual association과 더 비슷하고, unassociated hallucination이 더 쉽게 분리되었다는 것이다 \citep{cheang2026knowledge}.
분석에는 Meta-Llama-3-8B와 Mistral-7B-v0.3을 사용했고, 각 모델에서 12,293개 factual query를 분류했다.
\index{hallucination@환각(hallucination)}

분류 결과는 Llama에서 FA 3,506개, AH 1,406개, UH 7,381개였고, Mistral에서 각각 3,354개, 1,284개, 7,655개였다.
저자들은 subject representation patching, subject-to-last attention blocking과 last-token representation patching 뒤의 output distribution JS divergence를 layer별로 비교했다.
이 표본에서 FA와 AH의 intervention pattern은 유사했고 UH는 다른 pattern을 보였다.

내부 표현 probe도 같은 범위에서 AH보다 UH를 더 잘 검출했다.
예를 들어 Llama의 last-token representation probe AUROC는 AH-only에서 $0.69\pm0.03$, UH-only에서 $0.93\pm0.01$이었으며, Mistral에서는 각각 $0.63\pm0.02$와 $0.92\pm0.01$이었다.
이 수치는 모든 hallucination detector의 성능 상한이 아니라 논문이 정한 dataset, split, feature와 label에서의 결과다.

이 연구가 지지하는 제한된 해석은 내부 signal이 진실성 자체보다 parametric association 사용 여부를 더 강하게 반영할 수 있다는 것이다.
AH는 강하지만 잘못된 association을 통해 정답 회상과 비슷한 경로를 탈 수 있으므로, confidence나 hidden state만으로 참과 거짓을 안정적으로 가르기 어렵다는 설명이다.

\begin{limitbox}
probe 실패는 진실성 정보가 모델 내부 어디에도 존재하지 않는다는 증명이 아니다.
선택한 layer, position, 선형 feature와 dataset에서 충분히 복원되지 않았다는 결과다.
반대로 UH 검출 성공도 모델이 명시적인 ``나는 모른다'' 변수를 사용한다는 증명은 아니다.
\end{limitbox}

\section{in-context learning과 copying}

in-context learning(ICL)은 parameter update 없이 현재 context의 example이나 pattern에 따라 다음-token behavior가 달라지는 현상이다.
induction head는 반복 pattern의 복사라는 한 메커니즘을 제공하며, 소형 attention-only 모델에서는 인과적으로 자세히 분석되었다 \citep{olsson2022induction}.

일부 이론·실험 연구는 Transformer가 context 안의 linear regression example에서 gradient descent와 유사한 update를 구현할 수 있음을 보였다 \citep{vonoswald2023icl,akyurek2023icl}.
이 결과는 특정 synthetic regression task와 architecture에서 입출력 함수를 분석한 것이다.
자연어 ICL 전체가 실제 parameter gradient descent를 문자 그대로 실행한다는 결론은 아니다.

copying도 단일 기능이 아니다.
token을 그대로 반복하는 회로, 앞 token 다음의 token을 복사하는 induction, 이름 후보를 prediction position으로 옮기는 name mover와 긴 span을 재생하는 행동은 서로 다른 QK 조건과 OV 내용을 가질 수 있다.
``attention은 복사한다''는 문장은 source selection, value transformation과 target logit 효과를 분해해야 검증 가능하다.

\section{reasoning과 chain-of-thought}

chain-of-thought(CoT)는 모델이 중간 자연어 단계를 출력하도록 prompting하는 외부 생성 형식이다.
Wei 등은 충분히 큰 모델의 여러 arithmetic, commonsense와 symbolic reasoning benchmark에서 few-shot CoT prompt가 standard prompting보다 성능을 높인 결과를 보고했다 \citep{wei2022cot}.
이 결과는 출력된 문장이 내부 계산의 완전한 transcript임을 보이지 않는다.
\index{chain-of-thought@사고 연쇄(chain-of-thought)}

CoT의 비충실성은 입력의 biasing feature가 답을 바꾸면서도 생성된 explanation에는 그 영향이 언급되지 않는 실험으로 관찰되었다 \citep{turpin2023unfaithful}.
따라서 ``모델이 이유를 말했으므로 실제로 그 이유로 답했다''는 추론은 성립하지 않는다.
CoT는 행동 데이터이자 추가 계산을 유도하는 scratchpad가 될 수 있지만, 내부 메커니즘의 ground truth label은 아니다.

내부 reasoning 메커니즘을 주장하려면 중간 variable과 causal path를 별도로 시험해야 한다.
예를 들어 두 단계 질의에서 중간 entity direction을 찾고 이를 교체했을 때 최종 answer가 예측대로 바뀌는지 검사할 수 있다.
정답률 상승, probe와 intervention은 서로 다른 증거이므로 하나의 ``reasoning 능력'' 점수로 합치지 않는다.

\section{entity tracking, planning과 refusal}

entity tracking은 같은 이름이나 대상을 문맥의 여러 position에서 구분하고 해당 정보를 prediction position에 유지하는 계산을 요구한다.
IOI 회로는 duplicate token과 이름 억제, name-moving head가 협력하는 한 구체적 사례를 제공하지만, 임의 길이 대화의 모든 entity state를 설명하지는 않는다 \citep{wang2022ioi}.

2025년 Anthropic의 attribution-graph 사례 연구는 Claude 3.5 Haiku에서 multi-hop reasoning, 시의 운율 단어를 미리 고르는 planning, familiar entity와 unfamiliar entity의 구분, harmful request와 refusal에 관한 feature 경로를 보고했다 \citep{lindsey2025biology}.
저자들은 feature direction 개입과 후속 perturbation으로 선택한 그래프 가설을 검증했다.

이 결과는 frontier production model 내부에서 사람이 해석할 수 있는 중간 계산이 존재할 수 있다는 existence proof다.
그러나 저자들은 만족스러운 insight를 얻은 prompt가 시도한 prompt의 약 4분의 1이었고, 성공 사례에서도 모델 메커니즘의 작은 일부만 포착했으며, 선정한 사례가 도구의 한계에 의해 편향되었다고 명시했다.

거절에 관한 feature가 발견되었다고 모델의 안전 정책 전체가 한 방향에 저장된 것은 아니다.
pretraining과 post-training feature의 상호작용, prompt-specific graph와 여러 competing path가 존재할 수 있다.
발견된 방향을 steering했을 때 behavior가 바뀌는 결과도 원래 입력에서 그 방향이 모든 거절을 충분히 설명한다는 주장과 구분해야 한다.

\section{대형 모델의 attribution graph}

attribution graph 방법은 원 모델 자체를 곧바로 완전히 분해하지 않고, MLP를 더 해석하기 쉬운 cross-layer transcoder(CLT)로 바꾼 replacement model을 사용한다 \citep{ameisen2025circuittracing}.
CLT feature는 한 layer의 residual stream을 읽고 현재와 뒤 MLP layer의 출력을 재구성하며, sparsity penalty와 reconstruction loss로 학습된다.
\index{attribution graph@기여도 그래프(attribution graph)}

local replacement model은 한 prompt에서 원 attention pattern과 normalization denominator를 고정하고 CLT reconstruction error를 더하여 원 activation과 logit을 일치시킨다.
이 조정은 그 prompt의 forward value를 맞추지만, perturbation에 대한 반응까지 원 모델과 같다는 보장은 없다.
저자들은 이 차이를 mechanistic faithfulness로 따로 평가한다.

attribution graph의 edge는 고정된 비선형 조건 아래 feature 사이의 선형 효과를 나타낸다.
graph에는 token embedding, 활성 feature, error node와 output logit이 들어가며, 출력에 대한 기여가 작은 node와 edge를 pruning한다.
이 구조는 한 prompt를 읽기 쉬운 graph로 줄이지만, pruning threshold 아래의 다수 경로와 attention pattern을 만든 QK 메커니즘을 놓칠 수 있다.

replacement model의 근사 오차는 현재 방법의 핵심 한계다.
methods paper에서 가장 큰 18-layer CLT replacement model은 선택된 pretraining-style prompt 집합의 약 50\%에서 원 모델과 같은 top next token을 냈다.
local error correction으로 한 prompt의 수치를 맞춰도 동일한 내부 메커니즘을 사용한다는 결과가 자동으로 따라오지 않는다.

\begin{interpretbox}{현재 attribution graph의 올바른 지위}
attribution graph는 원 모델에 대한 구체적이고 검증 가능한 메커니즘 가설을 생성하는 강력한 현미경이다.
그것은 원 모델의 모든 계산을 손실 없이 사람이 읽을 수 있는 graph로 변환한 완전한 해부도가 아니다.
\end{interpretbox}

\section{모델 간 일반화 문제}

같은 decoder-only Transformer 계열도 factual recall의 중요한 구성요소가 같다고 가정할 수 없다.
Choe 등은 GPT, Llama, Qwen과 DeepSeek 계열을 비교하고, 조사한 Qwen-based 모델에서는 이전 GPT-style pattern과 달리 earliest layer attention module의 기여가 MLP보다 크다는 결과를 보고했다 \citep{choe2025crossarchitecture}.

모델 간 차이는 architecture 이름 하나로 설명되지 않는다.
normalization, GQA, tokenizer와 block 배치뿐 아니라 학습 데이터, objective와 seed가 다르다.
같은 기능을 다른 parameter 경로가 구현할 수도 있고, 표면적으로 비슷한 head가 다른 representation basis를 사용할 수도 있다.

일반화를 주장하려면 여러 수준을 분리해 측정해야 한다.
behavioral generalization은 같은 정답을 내는지 묻고, representational generalization은 alignment 뒤 유사한 feature가 있는지 묻고, mechanistic generalization은 같은 causal dependency가 intervention으로 유지되는지 묻는다.
첫째가 같아도 둘째와 셋째가 자동으로 같아지지 않는다.

2026년의 LM-agent 기반 circuit explanation 연구도 자동 설명의 범위를 명확히 보여준다.
HyVE는 84개 semi-synthetic Transformer circuit과 163개 component annotation으로 만든 benchmark에서 관찰, 가설, causal validation을 반복했고, Llama-3-8B arithmetic case study도 제시했다.
그러나 어떤 backbone도 일관되게 최선이 아니었고 validation 단계의 실패가 주요 장애로 남았다고 보고했다 \citep{khan2026agents}.

\section{Phase 7 학습 점검}

\begin{studybox}
\begin{enumerate}
  \item Geva et al. (2023)의 세 단계 각각에 대해 model, position, layer 범위, intervention, metric과 수치를 원문 표로 복원한다.
  \item 2510.09033의 FA/AH/UH labeling code를 재현하고 JS threshold가 바뀔 때 표본 구성이 얼마나 달라지는지 sensitivity curve를 그린다.
  \item factual recall 결과를 paraphrase, multi-token attribute와 다른 checkpoint에서 반복하고 실패 사례를 별도 표로 기록한다.
  \item CoT 문장과 내부 causal variable이 일치하는지 answer-preserving, rationale-changing counterfactual로 시험한다.
  \item attribution graph의 reconstruction error node, pruning threshold와 perturbation validation을 함께 보고한다.
  \item behavior, representation과 mechanism의 모델 간 일반화를 각각 다른 metric으로 평가한다.
\end{enumerate}
\end{studybox}

\begin{factbox}
Phase 7의 종료 조건은 유명한 회로 이름을 암기하는 것이 아니다.
각 논문의 결론을 원 실험 범위로 되돌리고, 관찰된 단계와 연구자가 붙인 해석을 구분하며, 새로운 모델에서 무엇을 다시 검증해야 하는지 설계할 수 있어야 한다.
\end{factbox}
